\emph{Pour chacune des affirmations suivantes, indiquer si elle est vraie ou    fausse.\\
	Justifier chaque réponse.\\
	Une réponse non justifiée ne rapporte aucun point.}

	\begin{enumerate}
		\item Le plan est muni d'un repère orthonormé.
		\begin{enumerate}
			\item On considère l'équation différentielle \quad $\left(E_{1}\right): \quad y'+2 y=0$,\quad où $y$ est une fonction de la variable $x$ définie sur $\R$.

			On note $f$ la solution de l'équation différentielle $(E_{1})$ telle que $f(0)=1$.

			\textbf{Affirmation 1 :} La courbe représentative de la fonction $f$ passe par le point A de coordonnées $\left(\ln 2 ; \dfrac{1}{2}\right)$.

			\item On considère la fonction $f$ définie sur l'intervalle $\left[\dfrac{\pi}{2} ; \pi\right]$ par :
			\[f(x)=x-\cos (x)\]

			\textbf{Affirmation 2 :} Sur l'intervalle $\left[\dfrac{\pi}{2} ; \pi\right]$, la courbe représentative de la fonction $f$ est au-dessus de ses tangentes.

			\item On considère la fonction $f$ définie sur l'intervalle $] 0 ;+\infty[$ par :
			\[f(x)=1+\dfrac{\sin (x)}{x}\]

			\textbf{Affirmation 3 :} La droite d'équation $y=1$ est une asymptote horizontale de la courbe représentative de la fonction $f$.
		\end{enumerate}

		\item On considère l'équation différentielle \quad $\left(E_{2}\right): \quad y'-3 y=2-6 x$,\quad où $y$ est une fonction de la variable $x$ définie sur $\R$.

		\textbf{Affirmation 4 :} La fonction $f$ définie sur $\R$ par
		\[f(x)=2 x-\e^{3 x}\]
		est une solution de l'équation $(E_{2})$.
	\end{enumerate}
